% !TeX root = ../main.tex
\section{Các Game Liên Quan}
\label{sec:related}

\subsection{Cờ Đam (Checkers / Draughts)}

Cờ Đam~\citep{schaeffer2007checkers} là game chiến thuật 2 người trên bàn
8{\texttimes}8. Quân di chuyển theo đường chéo và nhảy qua quân đối thủ
để ăn---quân bị ăn bị loại khỏi bàn. Tactical Formation giữ lại cơ chế
nhảy nhưng bỏ việc loại quân: quân bị nhảy qua trở thành thành phần
chiến lược thay vì bị xóa.

\subsection{Reversi / Othello}

Reversi là game 2 người trên bàn 8{\texttimes}8; chương trình Logistello
của \citet{buro1997othello} từng đánh bại nhà vô địch thế giới và đánh
dấu cột mốc AI cho game này. Gần đây, Othello $8\times8$ thậm chí đã được
chứng minh là \emph{hòa} với lối chơi tối ưu của cả hai bên%
~\citep{takizawa2023othello}, khép lại câu hỏi lý thuyết về game gốc.
Người chơi đặt quân mới và lật tất cả quân đối thủ bị kẹp theo đường thẳng. Tactical Formation giữ cơ chế lật quân
sandwich nhưng thay vì đặt quân mới, người chơi phải \emph{di chuyển quân
hiện có}. Ràng buộc này giới hạn đáng kể số nước khả dụng và thêm chiều
sâu về bố trí quân (positioning).

\subsection{Thiết Kế Game Cờ Trừu Tượng Mới}

Việc đề xuất một game cờ trừu tượng mới kèm phân tích bằng AI là một
dòng nghiên cứu đã có tiền lệ. Hệ thống \emph{Ludi} của
\citet{browne2010evolutionary} sinh tự động các luật game cờ trừu tượng
mới qua tiến hóa và đánh giá chất lượng game; một trong các kết quả của
nó---\emph{Yavalath}---kết hợp hai cơ chế cổ điển (xếp hàng để thắng và
phản-thắng) và đã được xuất bản thương mại~\citep{browne2012yavalath}.
Tuy nhiên, theo hiểu biết của chúng tôi, chưa có công trình nào lai cơ
chế \emph{nhảy} của Cờ Đam với cơ chế \emph{kẹp lật} của Reversi, cũng
như mở rộng lớp game Reversi-like sang nhiều người chơi---khoảng trống
mà bài báo này hướng tới.

\subsection{Các Game Di Chuyển Quân Khác}

Hai game cờ trừu tượng khác làm rõ vị trí của Tactical Formation trong
phổ thiết kế. \emph{Lines of Action}---game 2 người trên bàn
8{\texttimes}8 do Claude Soucie thiết kế, được mô tả trong tuyển tập của
\citet{sackson1969gamut}---cũng \emph{di chuyển quân hiện có} thay vì đặt
quân mới, và lấy việc \emph{kết nối bố trí} quân làm mục tiêu thắng; trò
chơi này đã được nghiên cứu sâu bằng tìm kiếm cây Monte
Carlo~\citep{winands2010loa}. Tuy nhiên, quân bị bắt trong Lines of Action
vẫn bị \emph{loại} khỏi bàn như ở Cờ Đam. \emph{Focus} (còn gọi
\emph{Domination}) của Sid Sackson~\citep{sackson1969gamut} lại gần với
Tactical Formation ở một điểm khác: quân bị bắt \emph{không bị xóa vĩnh
viễn} mà được giữ lại và có thể quay trở lại bàn cờ---một tiền lệ thiết kế
cho ý tưởng ``bắt quân là chuyển hóa, không phải loại bỏ''. Tactical
Formation hợp nhất hai nét này---di chuyển quân hiện có \emph{và} bảo toàn
quân số---rồi bổ sung cơ chế kẹp-lật kiểu Reversi và mở rộng lên tới 4
người chơi.

\subsection{So Sánh Tổng Quan}

\begin{table}[H]
  \centering
  \caption{So sánh Tactical Formation với các game cờ liên quan.}
  \label{tab:comparison}
  \begin{tabular}{lccc}
    \toprule
    \textbf{Đặc điểm} & \textbf{Cờ Đam} & \textbf{Reversi}
                      & \textbf{Tactical Formation} \\
    \midrule
    Cách di chuyển   & Nhảy qua quân  & Đặt quân mới & Nhảy qua quân \\
    Quân bị nhảy qua & Bị loại        & ---          & Còn trên bàn  \\
    Cách bắt quân    & Nhảy trực tiếp & Sandwich     & Sandwich sau nhảy \\
    Số người chơi    & 2              & 2            & 2, 3, hoặc 4  \\
    Số lớp vị trí    & 1              & 1            & 2 (đỉnh + tâm) \\
    \bottomrule
  \end{tabular}
\end{table}
